%% kapa: exam=1 round=28 year=2017 period=2 form=A subject=accounting subject_ko=회계학 seq=39 number=79
%% source: Q-Net 감정평가사 기출문제 (cst00309) artlSeq=5199603 file=extracted/1008/5199603/2017년 제28회 감정평가사 1차시험 문제지/2교시 A형.pdf
%% conversion: kordoc 4.24.51 -> tools/gen1.py (auto); review=proofread-sonnet; flags=none
\begin{question}{79}
(주)감평은 A제품을 생산ㆍ판매하고 있다. 20x1년에는 기존고객에게 9,000단위를 판매할 것으로 예상되며, A제품 관련 자료는 다음과 같다.
\par{\small\centering\begin{tabular}{|p{0.45\linewidth}|r|}
\hline
연간 최대생산량 & 10,000단위 \\ \hline
단위당 판매가격 & ￦2,000 \\ \hline
단위당 변동제조원가 & ￦1,000 \\ \hline
단위당 변동판매비 & ￦200 \\ \hline
연간 총고정제조원가 & ￦2,500,000 \\ \hline
\end{tabular}\par}
20x1년 중에 (주)감평은 새로운 고객인 (주)대한으로부터 A제품 2,000단위를 구매하겠다는 특별주문을 제안 받았다. 특별주문을 수락하면 기존고객에 대한 판매량 중 1,000단위를 감소시켜야 하며, 특별주문에 대해서는 단위당 변동판매비 ￦200이 발생하지 않는다. (주)감평이 특별주문으로부터 받아야 할 단위당 최소 판매가격은? (단, 특별주문은 일부분만 수락할 수 없음) \par
\begin{choices}
  \item ￦1,300
  \item ￦1,350
  \item ￦1,400
  \item ￦1,450
  \item ￦1,500
\end{choices}
\end{question}
\answer{3}
